%
% 第11章：数据集处理

\clearpage
\cleardoublepage

\chapter{数据集处理}

\section{数据收集}

任何机器学习系统的基础是高质量的数据。

\subsection{数据来源}

\begin{itemize}
    \item \textbf{公开数据集：} ImageNet、COCO、MNIST等
    \item \textbf{网页抓取：} 从网站自动收集数据
    \item \textbf{传感器：} 相机、麦克风、医疗设备
    \item \textbf{人工标注：} 专家标注
    \item \textbf{合成数据：} 从模拟生成
\end{itemize}

\begin{note}[数据质量]
垃圾输入，垃圾输出。高质量数据比复杂的算法更重要。
\end{note}

\section{数据标注}

标签使监督学习成为可能。

\subsection{标注类型}

\begin{itemize}
    \item \textbf{分类：} 每个图像单个标签（猫 vs 狗）
    \item \textbf{多标签：} 每个图像多个标签
    \item \textbf{目标检测：} 带标签的边界框
    \item \textbf{语义分割：} 像素级标签
    \item \textbf{实例分割：} 每个对象的实例ID
    \item \textbf{关键点：} 感兴趣点（人脸、姿态）
\end{itemize}

\section{数据划分}

正确的数据划分对无偏评估至关重要。

\subsection{基本划分}

\begin{itemize}
    \item \textbf{训练集：} 用于模型训练
    \item \textbf{验证集：} 用于超参数调优
    \item \textbf{测试集：} 用于最终评估（仅一次）
\end{itemize}

典型比例：
\begin{itemize}
    \item 小数据集（<10K）：70/10/20 或 80/10/10
    \item 大数据集（>1M）：98/1/1可能足够
\end{itemize}

\begin{note}[测试集泄露]
永远不要在训练期间使用测试集做任何决定。这会导致对测试集的过拟合和不可靠的性能估计。
\end{note}

\section{数据增强}

通过转换现有数据人为增加数据集大小。

\subsection{图像增强}

\begin{itemize}
    \item \textbf{旋转：} 随机角度旋转
    \item \textbf{翻转：} 水平/垂直翻转
    \item \textbf{裁剪：} 随机裁剪
    \item \textbf{颜色抖动：} 亮度、对比度、饱和度
    \item \textbf{缩放：} 随机缩放
    \item \textbf{平移：} 随机偏移
    \item \textbf{噪声：} 高斯噪声注入
    \item \textbf{Cutout：} 随机矩形遮挡
    \item \textbf{Mixup：} 混合两张图像和标签
    \item \textbf{CutMix：} 在图像之间剪切和粘贴补丁
\end{itemize}

\begin{code}
    \captionof{listing}{PyTorch中的数据增强}
    \label{code:augmentation_chinese}
\begin{minted}{python}
import torchvision.transforms as T

train_transform = T.Compose([
    T.RandomResizedCrop(224),
    T.RandomHorizontalFlip(p=0.5),
    T.RandomRotation(15),
    T.ColorJitter(brightness=0.2, contrast=0.2, 
                  saturation=0.2, hue=0.1),
    T.ToTensor(),
    T.Normalize(mean=[0.485, 0.456, 0.406],
                std=[0.229, 0.224, 0.225])
])

val_transform = T.Compose([
    T.Resize(256),
    T.CenterCrop(224),
    T.ToTensor(),
    T.Normalize(mean=[0.485, 0.456, 0.406],
                std=[0.229, 0.224, 0.225])
])
\end{minted}
\end{code}

\section{数据加载}

高效的数据加载对训练速度至关重要。

\begin{code}
    \captionof{listing}{PyTorch中的数据加载}
    \label{code:dataloader_chinese}
\begin{minted}{python}
from torch.utils.data import DataLoader, Dataset

class MyDataset(Dataset):
    def __init__(self, data, labels, transform=None):
        self.data = data
        self.labels = labels
        self.transform = transform
    
    def __len__(self):
        return len(self.data)
    
    def __getitem__(self, idx):
        image = self.data[idx]
        label = self.labels[idx]
        
        if self.transform:
            image = self.transform(image)
        
        return image, label

# 创建数据加载器
train_loader = DataLoader(
    train_dataset,
    batch_size=32,
    shuffle=True,  # SGD很重要
    num_workers=4,  # 并行数据加载
    pin_memory=True,  # 更快的GPU传输
    drop_last=True  # 一致的批次大小
)
\end{minted}
\end{code}

\section{归一化}

标准化输入特征以获得更好的训练。

图像归一化：
\begin{equation}
x_{norm} = \frac{x - \mu}{\sigma}
\end{equation}

常见ImageNet归一化：
\begin{itemize}
    \item 均值：$[0.485, 0.456, 0.406]$
    \item 标准差：$[0.229, 0.224, 0.225]$
\end{itemize}

\section{本章小结}

\begin{enumerate}
    \item 数据质量是模型性能的基础
    \item 标注类型取决于任务
    \item 正确的数据划分防止过拟合
    \item 数据增强改善泛化
    \item 高效数据加载加速训练
    \item 归一化稳定训练
\end{enumerate}
